\documentclass[11pt]{article}
\usepackage[a4paper,margin=1in]{geometry}
\usepackage{amsmath,amssymb,mathtools,bm}
\usepackage{enumitem,booktabs,array}
\usepackage{tcolorbox}
\usepackage{hyperref}
\usepackage[protrusion=true,expansion=false]{microtype}
\hypersetup{colorlinks=true,linkcolor=blue,urlcolor=blue,citecolor=blue}
\setlist[itemize]{topsep=2pt,itemsep=2pt}
\setlist[enumerate]{topsep=2pt,itemsep=2pt}
\newtcolorbox{keybox}{colback=blue!3!white,colframe=blue!40!black,boxrule=0.4pt,arc=1mm}
\newtcolorbox{alertbox}{colback=red!3!white,colframe=red!40!black,boxrule=0.4pt,arc=1mm}
\newtcolorbox{answerbox}{colback=green!3!white,colframe=green!40!black,boxrule=0.4pt,arc=1mm}
\newtcolorbox{drillbox}{colback=yellow!5!white,colframe=orange!60!black,boxrule=0.4pt,arc=1mm}
\newcommand{\EV}{\mathbb{E}}
\newcommand{\Var}{\mathrm{Var}}
\newcommand{\Cov}{\mathrm{Cov}}
\newcommand{\blank}[1]{\underline{\hspace{#1}}}

\title{E07-02: Bernstein \& Sheen (2016, RFS) \\
\large ``The Operational Consequences of Private Equity Buyouts: Evidence from the Restaurant Industry'' --- Active Recall Workbook}
\author{Empirical Corporate Finance II --- Exam Prep}
\date{\today}
\begin{document}\maketitle

\begin{keybox}
\textbf{Paper in one sentence.} BS use health-inspection data as an objective operational-quality measure for $\sim$50 restaurant-chain PE buyouts in Florida and show that PE-owned chains meaningfully \emph{improve} health, safety, and maintenance standards --- providing detailed evidence that PE does create operational value at the outlet level, at least in the restaurant industry.

\textbf{Placement.} Complements BDS (2014): while plant-level TFP gains are elusive, micro-level operational improvements can be real in industries with measurable service-quality proxies.
\end{keybox}

\section*{Round 1: Complete Walk-Through}

\subsection*{1.1 Research design}
Use a setting where operational quality is directly measured: restaurant health-and-safety inspections. Compare PE-owned vs.\ non-PE-owned outlets within the same chain and/or versus matched controls.

\subsection*{1.2 Data}
\begin{itemize}
\item Florida Department of Business and Professional Regulation: $\sim$500,000 inspection records 2002--2012.
\item $\sim$50 PE-backed restaurant-chain buyouts.
\item Outlet-level violation counts: critical (food safety), non-critical (maintenance), total.
\end{itemize}

\subsection*{1.3 Empirical design}
DiD using outlet-level panel:
\[
\text{Violations}_{oct}=\alpha_o+\delta_t+\beta\cdot\text{Post-PE}_{oct}+X_{oct}'\gamma+\varepsilon_{oct},
\]
where $o$ = outlet, $c$ = chain, $t$ = time. Outlet and time fixed effects soak up persistent quality differences.

\subsection*{1.4 Headline results}
\begin{itemize}
\item Post-PE, outlets show 5--15\% fewer critical violations, 10--20\% fewer non-critical.
\item Effect strongest for chains that invest in centralized training, supply-chain, and maintenance systems post-buyout.
\item Menu prices and employment also shift: some outlets lay off workers, raise prices modestly.
\item Bankruptcies and closures slightly more likely at low-performing pre-PE outlets.
\end{itemize}

\subsection*{1.5 Mechanism}
\begin{itemize}
\item PE owners impose standardized operating procedures across outlets.
\item Investment in IT, supply chain, and training raises minimum standards.
\item Concentrated ownership internalizes brand externalities from local-outlet quality.
\end{itemize}

\subsection*{1.6 Implications}
\begin{itemize}
\item PE effects are real in service industries with observable operational proxies.
\item Challenges the view (BDS) that PE ownership has no real-economy impact.
\item Highlights value of micro outcomes (inspections, customer reviews) over aggregate TFP.
\end{itemize}

\subsection*{1.7 Caveats}
\begin{alertbox}
\begin{itemize}
\item Florida-restaurant sample limits generalizability.
\item PE-targeted chains may differ on unobservables.
\item Inspection-recording practices may change post-acquisition.
\end{itemize}
\end{alertbox}

\section*{Round 2: Level 1 Fill-in}
\begin{enumerate}
\item Data: Florida \blank{2cm} inspections.
\item Sample: \blank{1cm} PE restaurant-chain buyouts.
\item Outcome: post-PE critical violations fall \blank{1cm}--\blank{1cm}\%.
\item Mechanism: centralized \blank{2cm} and supply chain.
\item Contrast with \blank{1cm} (2014) plant-level null.
\end{enumerate}

\section*{Round 3: Level 2 Prompts}
\begin{enumerate}
\item Describe the health-inspection data and why it is a good proxy.
\item Report the main findings on violations, employment, and prices.
\item Discuss the mechanisms PE operators use.
\item Compare BS with Bharath-Dittmar-Sivadasan (2014).
\item Evaluate external validity of restaurant-industry findings.
\end{enumerate}

\section*{Round 4: Minimal Scaffolding}
From a blank page: (i) setting, (ii) data, (iii) results, (iv) mechanism, (v) context.

\section*{Round 5: Blank Page}
Essay: PE ownership and operational quality --- evidence from restaurants.

\vspace{6pt}
\begin{drillbox}
\textbf{Speed Drill.} (1) Data. (2) Sample size. (3) Violation change. (4) Mechanism. (5) Contrast paper.
\end{drillbox}

\section*{Quick Reference \& Answer Key}
\begin{answerbox}
\begin{itemize}
\item \textbf{Data.} FL health inspections.
\item \textbf{Result.} Violations drop 5--20\% post-PE.
\item \textbf{Mechanism.} Standardization, training, supply chain.
\item \textbf{Lesson.} PE real ops effects do exist in some settings.
\end{itemize}
\end{answerbox}

\begin{keybox}
\textbf{30-second exam pitch.} Bernstein \& Sheen (2016) use Florida restaurant-chain health inspections as an objective measure of operational quality to test whether PE ownership creates real operational value. With $\sim$500,000 inspection records 2002--2012 covering $\sim$50 PE restaurant-chain buyouts in an outlet-level DiD, they find that post-acquisition, outlets show 5--15\% fewer critical violations and 10--20\% fewer non-critical violations, with effects strongest in chains investing in centralized training, supply chains, and maintenance. Some outlets see modest price increases and employment shifts; low-performing outlets are more likely to close. Mechanism: PE owners impose standardized operating procedures and internalize brand externalities across outlets. The paper complements Bharath-Dittmar-Sivadasan (2014) --- where plant-level TFP gains were elusive --- by showing PE does create micro-operational value in industries with observable service-quality proxies, and motivates richer outcome measures in PE research.
\end{keybox}

\end{document}
